% GENERATED FILE. Edit canonical lessons, then run npm run generate:books.
% canonical-source-hash: fnv1a64:4ef08fa8aa5120bd
% canonical-lessons: ZH-C310-huaiyi, ZH-C310-zhengren, ZH-C310-zhengju, ZH-C310-shuohuang, ZH-C310-chengren

\chapter{Who Did It}
\label{ch:zh-who-did-it}
\hlchaptermodality{310}

\begin{chapteropening}
\textbf{By the end of this chapter:} \emph{I can talk about a suspect, a witness and evidence, lying and owning up.}

Complete the last lesson of chapter 310: I can talk about a suspect, a witness and evidence, lying and owning up.
\end{chapteropening}

\section[huáiyí]{huáiyí — to suspect; to doubt}

\label{lesson:ZH-C310-huaiyi}



\begin{quote}
Before the new one: say the Mandarin for a con man, then the Mandarin for on purpose.
\end{quote}

\subsection*{You'll want to know: huáiyí}
\emph{huáiyí} — \textquotedblleft{}to suspect; to doubt\textquotedblright{}.

\begin{quote}
\emph{jǐngchá huáiyí tā shì xiǎotōu} — The police suspect he's the thief.
\end{quote}

Its characters come later; for now it is a word for the ear.

\begin{cousinweb}
\emph{huáiyí} is built from \emph{huái} 'to hold close' and \emph{yí} 'to doubt', both inherited from Middle Chinese, \emph{hweaj} and \emph{ngi} (Guangyun, in Baxter's transcription), which glosses the first 'to embrace' and 'to think of' and the second 'unsure' and 'to suspect': to hold a doubt.
\end{cousinweb}

\begin{grammarlens}[title={What you've built}]
One more word for telling what happened.
\end{grammarlens}

\subsection*{Your turn}
\begin{itemize}
\raggedright
  \item \emph{Say it:} \emph{huáiyí}
  \item \emph{Say it:} \emph{huáiyí}, once more
  \item \emph{Say it:} \emph{jǐngchá huáiyí tā shì xiǎotōu} — The police suspect he's the thief.
  \item \emph{Recall:} say the Mandarin for a con man, then the Mandarin for on purpose, then say \emph{huáiyí} again
\end{itemize}

\begin{tcolorbox}[breakable,colback=teal!4,colframe=teal!35!black,title={Before you move on}]
What does huáiyí mean? (\textquotedblleft{}To suspect; to doubt\textquotedblright{}.) Say it once more.
\end{tcolorbox}
\section[zhèngrén]{zhèngrén — witness}

\label{lesson:ZH-C310-zhengren}



\begin{quote}
Before the new one: say the Mandarin for on purpose, then the Mandarin for to suspect someone.
\end{quote}

\subsection*{You'll want to know: zhèngrén}
\emph{zhèngrén} — \textquotedblleft{}witness\textquotedblright{}.

\begin{quote}
\emph{tā shì zhège shìgù de zhèngrén} — He's a witness to the accident.
\end{quote}

Its characters come later; for now it is a word for the ear.

\begin{cousinweb}
\emph{zhèngrén} is built from \emph{zhèng} 'proof' and \zh{人} 'a person', both inherited from Middle Chinese, \emph{tsyingH} and \emph{nyin} (Guangyun, in Baxter's transcription), which glosses the first 'to verify'.
\end{cousinweb}

\begin{grammarlens}[title={What you've built}]
One more word for telling what happened.
\end{grammarlens}

\subsection*{Your turn}
\begin{itemize}
\raggedright
  \item \emph{Say it:} \emph{zhèngrén}
  \item \emph{Say it:} \emph{zhèngrén}, once more
  \item \emph{Say it:} \emph{tā shì zhège shìgù de zhèngrén} — He's a witness to the accident.
  \item \emph{Recall:} say the Mandarin for on purpose, then the Mandarin for to suspect someone, then say \emph{zhèngrén} again
\end{itemize}

\begin{tcolorbox}[breakable,colback=teal!4,colframe=teal!35!black,title={Before you move on}]
What does zhèngrén mean? (\textquotedblleft{}Witness\textquotedblright{}.) Say it once more.
\end{tcolorbox}
\section[zhèngjù]{zhèngjù — evidence; proof}

\label{lesson:ZH-C310-zhengju}



\begin{quote}
Before the new one: say the Mandarin for to suspect someone, then the Mandarin for a witness to a crime.
\end{quote}

\subsection*{You'll want to know: zhèngjù}
\emph{zhèngjù} — \textquotedblleft{}evidence; proof\textquotedblright{}.

\begin{quote}
\emph{jǐngchá zhǎodào le zhèngjù} — The police found the evidence.
\end{quote}

Its characters come later; for now it is a word for the ear.

\begin{cousinweb}
\emph{zhèngjù} is built from \emph{zhèng} 'to verify' and \emph{jù} 'to lean on', both inherited from Middle Chinese, \emph{tsyingH} and \emph{kjoH} (Guangyun, in Baxter's transcription), which glosses the second 'to rely on' and 'to hold': what a proof can lean on.
\end{cousinweb}

\begin{grammarlens}[title={What you've built}]
One more word for telling what happened.
\end{grammarlens}

\subsection*{Your turn}
\begin{itemize}
\raggedright
  \item \emph{Say it:} \emph{zhèngjù}
  \item \emph{Say it:} \emph{zhèngjù}, once more
  \item \emph{Say it:} \emph{jǐngchá zhǎodào le zhèngjù} — The police found the evidence.
  \item \emph{Recall:} say the Mandarin for to suspect someone, then the Mandarin for a witness to a crime, then say \emph{zhèngjù} again
\end{itemize}

\begin{tcolorbox}[breakable,colback=teal!4,colframe=teal!35!black,title={Before you move on}]
What does zhèngjù mean? (\textquotedblleft{}Evidence; proof\textquotedblright{}.) Say it once more.
\end{tcolorbox}
\section[shuōhuǎng]{shuōhuǎng — to lie; to tell a lie}

\label{lesson:ZH-C310-shuohuang}



\begin{quote}
Before the new one: say the Mandarin for a witness to a crime, then the Mandarin for proof of what happened.
\end{quote}

\subsection*{You'll want to know: shuōhuǎng}
\emph{shuōhuǎng} — \textquotedblleft{}to lie; to tell a lie\textquotedblright{}.

\begin{quote}
\emph{wǒ juéde tā zài shuōhuǎng} — I think he's lying.
\end{quote}

Its characters come later; for now it is a word for the ear.

\begin{cousinweb}
\emph{shuōhuǎng} is built from \zh{说} 'to say' and \emph{huǎng} 'a lie'. The first is inherited from Middle Chinese \emph{sywet} (Guangyun, in Baxter's transcription); the Guangyun has no entry for the second, so its history is not settled.
\end{cousinweb}

\begin{grammarlens}[title={What you've built}]
One more word for telling what happened.
\end{grammarlens}

\subsection*{Your turn}
\begin{itemize}
\raggedright
  \item \emph{Say it:} \emph{shuōhuǎng}
  \item \emph{Say it:} \emph{shuōhuǎng}, once more
  \item \emph{Say it:} \emph{wǒ juéde tā zài shuōhuǎng} — I think he's lying.
  \item \emph{Recall:} say the Mandarin for a witness to a crime, then the Mandarin for proof of what happened, then say \emph{shuōhuǎng} again
\end{itemize}

\begin{tcolorbox}[breakable,colback=teal!4,colframe=teal!35!black,title={Before you move on}]
What does shuōhuǎng mean? (\textquotedblleft{}To lie; to tell a lie\textquotedblright{}.) Say it once more.
\end{tcolorbox}
\section[chéngrèn]{chéngrèn — to admit; to own up}

\label{lesson:ZH-C310-chengren}



\begin{quote}
Before the new one: say the Mandarin for proof of what happened, then the Mandarin for to tell a lie.
\end{quote}

\subsection*{You'll want to know: chéngrèn}
\emph{chéngrèn} — \textquotedblleft{}to admit; to own up\textquotedblright{}.

\begin{quote}
\emph{zuìhòu tā chéngrèn le} — In the end he owned up.
\end{quote}

Its characters come later; for now it is a word for the ear.

\begin{cousinweb}
\emph{chéngrèn} is built from \emph{chéng} 'to receive' and \emph{rèn} 'to recognize', both inherited from Middle Chinese, \emph{dzying} and \emph{nyinH} (Guangyun, in Baxter's transcription), which glosses the first 'to receive' and the second 'to know'.
\end{cousinweb}

\begin{grammarlens}[title={What you've built}]
One more word for telling what happened.
\end{grammarlens}

\subsection*{Your turn}
\begin{itemize}
\raggedright
  \item \emph{Say it:} \emph{chéngrèn}
  \item \emph{Say it:} \emph{chéngrèn}, once more
  \item \emph{Say it:} \emph{zuìhòu tā chéngrèn le} — In the end he owned up.
  \item \emph{Recall:} say the Mandarin for proof of what happened, then the Mandarin for to tell a lie, then say \emph{chéngrèn} again
\end{itemize}

\begin{tcolorbox}[breakable,colback=teal!4,colframe=teal!35!black,title={Before you move on}]
What does chéngrèn mean? (\textquotedblleft{}To admit; to own up\textquotedblright{}.) Say it once more.
\end{tcolorbox}
